\section{Exact Gaussian projection}\label{h:projection}

A block of exact equations acts on a localized part of each
continuation's successful mass. We prove its projection bound for an
arbitrary finite measure with controlled mass in every ambient ball.
Keeping the coefficient and support radius visible will expose the
cancellation in the block estimate. The same theorem will later apply
to a prior supported on a lower-dimensional part of the sphere.

For a positive integer \(m\), let \(\gamma_m\) be the law of an
\(m\)-by-\(d\) matrix with independent standard Gaussian entries, and set
\[
 \lambda_m(dA,dy)=\gamma_m(dA)\,dy.
\]
For a finite nonnegative Borel measure \(\mu\), let \(\eta_\mu\) be
the image of \(\gamma_m(dA)\mu(ds)\) under
\((A,s)\mapsto(A,As)\). The matrix remains part of the observed data;
no normalization of \(\mu\) is made. The integer \(m\) counts rows,
while \(p\) below is a density moment order. Write \(v_m\) for the unit
\(m\)-ball volume. For a closed support ball \(B(z,R)\), define
\begin{equation}\label{h:dv:support-sets}
 \mathcal E_{z,R}=\{(A,y):y\in Az+A B(0,R)\}
 \subseteq
 \mathcal S_{z,R}
 =\{(A,y):\|y-Az\|\le R\|A\|_{\mathrm{op}}\}.
\end{equation}
The exact image is carried by \(\mathcal E_{z,R}\). This set is
Borel: it is the zero set of the continuous function
\((A,y)\mapsto\min_{\|v\|\le R}\|y-Az-Av\|\).
Changing \((A,y)\) changes this minimum by at most
\(\|y-y'\|+\|A-A'\|_{\mathrm{op}}(\|z\|+R)\).

In this section \(H\) is a local-mass coefficient and \(R\) is the
geometric radius of a supporting ball. This use of \(R\) is separate
from the propagated scale in \eqref{h:local:bound}.

\begin{theorem}[Exact projection density from an ambient ball bound]
\label{h:dv:finite-density}
Let \(d,m\ge1\) and \(p\ge2\) be integers, and let
\(\beta>m+p-2\), \(H,R>0\), and \(z\in\mathbb R^d\).
Suppose that \(\mu\) is a finite nonnegative Borel measure supported
on \(B(z,R)\) and satisfying
\begin{equation}\label{h:dv:ambient-bound}
 \mu(B(x,r))\le Hr^\beta
 \qquad(x\in\mathbb R^d,\ r>0).
\end{equation}
For \(0\le j\le p-2\), put
\[
 J_j=1+3^j2^\beta\frac{m}{\beta-j-m},
\]
and define
\begin{equation}\label{h:dv:moment-scale}
 \mathcal B=(2\pi)^{-m(p-1)/(2p)}
 H R^{\beta-m(1-1/p)}
 \left(\prod_{j=0}^{p-2}J_j\right)^{1/p}.
\end{equation}
Then \(\eta_\mu\) has a nonnegative jointly Borel density \(F\)
with respect to \(\lambda_m\), which can be chosen to vanish
outside \(\mathcal E_{z,R}\), and
\begin{equation}\label{h:dv:finite-density-bound}
 \|F\|_{L^p(\lambda_m)}\le\mathcal B.
\end{equation}
\end{theorem}

Projection loses \(m(1-1/p)\) powers of the support radius.
The volume of the projected ball will restore exactly this power.
We prove the theorem by controlling inverse affine distances,
expanding a smoothed density moment, and then removing the smoothing
on matrix--label space. The independent-copy moment and successive
orthogonalization have a close predecessor in the noisy regression
analysis of \cite[Section~7, Lemma~12]{SSV2019}. That result treats a
one-row noisy interval under a global \(L^2\) density condition;
here the input is local growth, and the output concerns several rows
and exact labels.

\subsection{Affine distances and the Gram determinant}

\begin{proposition}[Affine-distance and Gram bounds]
\label{h:dv:ambient-gram}
Under the assumptions of Theorem~\ref{h:dv:finite-density}, use
the factors \(J_j\) defined there.
If \(L\) is an affine \(j\)-plane with \(0\le j\le p-2\), then,
for \(0<t\le R\),
\begin{align}
 \mu\{s:\operatorname{dist}(s,L)\le t\}
 &\le 3^j2^\beta H R^j t^{\beta-j},
 \label{h:dv:tube}\\
 \int\operatorname{dist}(s,L)^{-m}\,d\mu(s)
 &\le J_j H R^{\beta-m}.
 \label{h:dv:inverse-distance}
\end{align}
For a tuple \((s_1,\ldots,s_p)\), let
\(V=[s_2-s_1,\ldots,s_p-s_1]\).  The matrix \(V\) has full column rank
for \(\mu^{\otimes p}\)-almost every tuple, and
\begin{equation}
 \int\det(V^{\mathsf T}V)^{-m/2}\,d\mu^{\otimes p}
 \le
 H^p R^{\beta p-m(p-1)}
 \prod_{j=0}^{p-2}J_j.
 \label{h:dv:gram-product}
\end{equation}
The integrand in \eqref{h:dv:gram-product} may be assigned any value on
the null set of rank-deficient tuples.

\end{proposition}

\begin{proof}
The local growth hypothesis is of Frostman type, and
inverse-distance averaging is a classical projection method; compare
\cite[Lemmas~3.5 and~5.1, pp.~230--231 and~236]{Mattila1975}.
Here we retain the coefficient and radius throughout the calculation.

The orthogonal projection of \(B(z,R)\) onto \(L\) lies in a
\(j\)-dimensional ball of radius \(R\).  A maximal \(t\)-separated set
in that ball is a \(t\)-net and has at most
\((1+2R/t)^j\le(3R/t)^j\) points, by comparing disjoint balls of radius
\(t/2\) with the ball of radius \(R+t/2\).  When \(j=0\), use its
single point.  Every point of \(B(z,R)\) at distance at most \(t\) from
\(L\) is within \(2t\) of a net point.  Applying
\eqref{h:dv:ambient-bound} to these ambient balls proves
\eqref{h:dv:tube}.  Since \(\beta-j>0\), letting \(t\downarrow0\)
also gives \(\mu(L)=0\).

The total mass is at most \(HR^\beta\).  Integrating the elementary
identity for an inverse power up to radius \(R\), and using Tonelli's
theorem, gives
\begin{align*}
 \int\operatorname{dist}(s,L)^{-m}\,d\mu(s)
 &\le R^{-m}\mu(\mathbb R^d)
   +m\int_0^R t^{-m-1}
       \mu\{s:\operatorname{dist}(s,L)\le t\}\,dt\\
 &\le
 \left(1+3^j2^\beta\frac{m}{\beta-j-m}\right)
 HR^{\beta-m}.
\end{align*}
The denominator is positive by the stated hypothesis on \(\beta\).
This proves \eqref{h:dv:inverse-distance}.

If \(p-1>d\), then \(d\le p-2\); applying the tube estimate to
\(L=\mathbb R^d\) and letting \(t\downarrow0\) shows that \(\mu=0\).
The Gram conclusion is then immediate.  Otherwise, the nullity of
affine planes of dimensions at most \(p-2\) shows by successive
integration that \(s_i\) lies outside
\(\operatorname{aff}(s_1,\ldots,s_{i-1})\) for almost every tuple and
every \(2\le i\le p\).  Hence \(V\) has full column rank almost
everywhere.  Gram--Schmidt orthogonalization gives, on this set,
\[
 \det(V^{\mathsf T}V)^{-m/2}
 =
 \prod_{i=2}^p
 \operatorname{dist}\bigl(s_i,
       \operatorname{aff}(s_1,\ldots,s_{i-1})\bigr)^{-m}.
\]
The product of the corresponding positive heights is
\((p-1)!\) times the \((p-1)\)-simplex volume. Related
simplex-volume factors occur in Drury's affine-plane Jacobian
\cite[Lemma~1, pp.~497--498]{Drury1984}; its normalization and
exponent differ from this Gaussian calculation.

Integrate first in \(s_p\), then in \(s_{p-1}\), and continue down to
\(s_2\).  At the \(s_i\) step the preceding affine span has dimension
\(i-2\) almost everywhere, so \eqref{h:dv:inverse-distance} contributes
\(J_{i-2}HR^{\beta-m}\).  The final \(s_1\) integral contributes at
most \(HR^\beta\).  All integrands are nonnegative, and Tonelli's
theorem therefore justifies this order and proves
\eqref{h:dv:gram-product}.

\end{proof}

\subsection{Smoothing and exact labels}

Let \(\phi_\delta\) be the \(N(0,\delta^2I_m)\) density on
\(\mathbb R^m\). We first record the Gaussian identity for all tuples,
including rank-deficient ones.

\begin{lemma}[Gaussian difference identity]\label{h:proj:gaussian}
Let \(d,m\ge1\) and \(k\ge2\) be integers, let \(\delta>0\), and
let \(t_1,\ldots,t_k\in\R^d\). Put
\(D=[t_2-t_1,\ldots,t_k-t_1]\). If \(A\) has law \(\gamma_m\), then
\begin{equation}\label{h:proj:gaussian-identity}
\begin{split}
 \E_A\int_{\R^m}\prod_{i=1}^k\phi_\delta(y-At_i)\,dy
 ={}&(2\pi)^{-m(k-1)/2}\\
 &{}\times
 \det\!\left(D^{\mathsf T}D+
 \delta^2(I_{k-1}+\mathbf1\mathbf1^{\mathsf T})\right)^{-m/2}.
\end{split}
\end{equation}
If \(D\) has full column rank, the right side is at most
\((2\pi)^{-m(k-1)/2}\det(D^{\mathsf T}D)^{-m/2}\).
\end{lemma}

\begin{proof}
Let \(\eta_1,\ldots,\eta_k\) be independent
\(N(0,\delta^2I_m)\) vectors. For fixed \(A\), integrating the
product over the common location \(y\) gives the density at zero of
the \(k-1\) differences
\[
 (At_i+\eta_i)-(At_1+\eta_1),\qquad 2\le i\le k.
\]
Indeed, the linear change from \(k\) vectors to the first vector and
these differences has determinant one. After averaging over \(A\),
each of the \(m\) rows is a centered Gaussian difference vector with
covariance
\[
 D^{\mathsf T}D+\delta^2(I_{k-1}+\mathbf1\mathbf1^{\mathsf T}).
\]
The \(m\) row vectors are independent. Evaluating their joint density
at zero gives \eqref{h:proj:gaussian-identity}. For fixed
\(\delta>0\), the conditional difference densities are continuous
and bounded uniformly in \(A\), so their averaged value at zero
equals the value of the averaged density. The covariance is positive
definite even when \(D\) is rank deficient.

If \(G=D^{\mathsf T}D\) is positive definite and \(B\) is positive
semidefinite, then
\(\det(G+B)=\det G\,\det(I+G^{-1/2}BG^{-1/2})\ge\det G\).
Apply this with \(B=\delta^2(I+\mathbf1\mathbf1^{\mathsf T})\).
\end{proof}

\Needspace{4\baselineskip}
\begin{lemma}[Gaussian smoothing with the full mass factor]
\label{h:dv:gaussian-moment}
Under the assumptions of Theorem~\ref{h:dv:finite-density}, use
the scale \(\mathcal B\) from \eqref{h:dv:moment-scale}.
Let \(\phi_\delta\) be the \(N(0,\delta^2I_m)\) density and put
\[
 F_\delta(A,y)=\int\phi_\delta(y-As)\,d\mu(s),
 \qquad \delta>0.
\]
Then \(F_\delta\) is jointly Borel, and, with
\(V=[s_2-s_1,\ldots,s_p-s_1]\),
\begin{align}
 \int F_\delta^p\,d\lambda_m
 &=(2\pi)^{-m(p-1)/2}
   \int\det\!\left(
     V^{\mathsf T}V+
     \delta^2(I_{p-1}+\mathbf 1\mathbf 1^{\mathsf T})
   \right)^{-m/2}\,d\mu^{\otimes p}
 \notag\\
 &\le \mathcal B^p.
 \label{h:dv:gaussian-moment-bound}
\end{align}
\end{lemma}

\begin{proof}
The defining integrand is nonnegative and Borel, so its parameter
integral is Borel.  Expand the \(p\)th power and use Tonelli's theorem.
The general Gaussian difference identity
\eqref{h:proj:gaussian-identity}, with \(k=p\), \(m\) rows, and the
tuple \(s_1,\ldots,s_p\), gives the displayed equality.  In particular
the common first noise contributes
\(\delta^2\mathbf1\mathbf1^{\mathsf T}\), and the prefactor is
\((2\pi)^{-m(p-1)/2}\).

On the full-rank tuples of Proposition~\ref{h:dv:ambient-gram},
adding the positive definite noise covariance can only increase the
determinant.  Indeed, for positive definite \(G\) and positive
semidefinite \(Q\),
\(\det(G+Q)=\det(G)\det(I+G^{-1/2}QG^{-1/2})\ge\det(G)\).
The rank-deficient tuples have product measure zero.  The bound now
follows from \eqref{h:dv:gram-product} and the definition of
\(\mathcal B\).
\end{proof}

To pass from the uniform moment bound to exact observations, we use
the following compact-test form of \(L^p\) duality.

\begin{lemma}[A density from compactly supported tests]
\label{h:proj:density-by-tests}
Let \(E\) be a Euclidean space, let \(\lambda\) be a locally finite,
sigma-finite Borel measure on \(E\), and let \(\eta\) be a finite
nonnegative Borel measure on \(E\). Suppose \(p>1\), \(U\ge0\), and
\[
 \left|\int\psi\,d\eta\right|
 \le U\|\psi\|_{L^{p/(p-1)}(\lambda)}
 \qquad\text{for every real }\psi\in C_c(E).
\]
Then \(\eta=f\lambda\) for a nonnegative
\(f\in L^p(\lambda)\) with \(\|f\|_p\le U\).
\end{lemma}

\begin{proof}
The locally finite Borel measures here have Radon extensions
\cite[Corollary~2.11]{StanfordRadonMeasures}; on Euclidean space
these are the completions of the Borel measures. On the completion of
\(\lambda\), the space \(C_c(E)\) is dense in
\(L^{p/(p-1)}(\lambda)\)
\cite[Proposition~416I]{FremlinRadonMeasures}. The assumed functional
extends by continuity to that space. Real \(L^p\) duality
\cite[Theorem~244K]{FremlinFunctionSpaces} gives a representing
\(f\in L^p(\lambda)\) of norm at most \(U\).

Choose a Borel representative of \(f\)
\cite[Section~241B(k)]{FremlinFunctionSpaces}. H\"older makes
\(f^+\lambda\) and \(f^-\lambda\) finite on compact sets. Their
Radon extensions obey
\[
 \int\psi\,d(f^+\lambda)
 =\int\psi\,d(\eta+f^-\lambda)
 \qquad(\psi\in C_c(E)).
\]
Radon measures are determined by these tests
\cite[Proposition~416E(b)]{FremlinRadonMeasures}. The two measures
therefore agree. Evaluating on the Borel set where \(f<0\) gives
\(f^-\lambda=0\), and then \(\eta=f\lambda\).
\end{proof}

\begin{proof}[Proof of Theorem~\ref{h:dv:finite-density}]
For a real \(\psi\in C_c(\mathbb R^{md}\times\mathbb R^m)\), let
\(Z\sim N(0,I_m)\) be independent.  Then
\[
 \int\psi F_\delta\,d\lambda_m
 =\int\mathbb E_Z\psi(A,As+\delta Z)\,
          d\gamma_m(A)d\mu(s)
 \longrightarrow\int\psi\,d\eta_\mu
\]
by bounded convergence, since \(\mu\) is finite.  H\"older's inequality
and Lemma~\ref{h:dv:gaussian-moment} bound the limiting functional by
\(\mathcal B\|\psi\|_{L^{p/(p-1)}(\lambda_m)}\).
The measure \(\lambda_m\) is locally finite and \(\sigma\)-finite
Borel, and \(\eta_\mu\) is finite nonnegative Borel.  Applying
Lemma~\ref{h:proj:density-by-tests} with \(U=\mathcal B\) gives
\eqref{h:dv:finite-density-bound}. Choose a Borel representative
and multiply it by the indicator of \(\mathcal E_{z,R}\).
The image measure is supported on
the Borel set \(\mathcal E_{z,R}\), so its nonnegative density vanishes
almost everywhere outside that set.
\end{proof}

\Needspace{4\baselineskip}
\begin{corollary}[Normalized exact projection]\label{h:proj:normalized}
Use the parameters in \eqref{h:local:parameters}, and let
\(1\le m\le q\). Let \(\nu\) be a finite nonnegative Borel measure
of mass at most one, supported on the unit ball in \(\mathbb R^d\),
such that
\begin{equation}\label{h:proj:growth}
 \nu(B(t,\rho))\le\rho^a
 \qquad(t\in\supp\nu,\ \rho>0).
\end{equation}
The image of \(\gamma_m(dA)\nu(dt)\) under
\((A,t)\mapsto(A,At)\) has a nonnegative jointly Borel density
\(f_\nu\) relative to \(\lambda_m\), and
\begin{equation}\label{h:proj:norm}
 \|f_\nu\|_{L^q(\lambda_m)}\le\exp(Cd).
\end{equation}
The constant is independent of \(\nu\) and \(m\).
\end{corollary}

\begin{proof}
The zero measure is immediate. Any ambient ball of radius \(\rho\)
with positive mass contains a support point, and is contained in
the ball of radius \(2\rho\) about that point. Hence its mass is at
most \((2\rho)^a\). Apply Theorem~\ref{h:dv:finite-density} with
\(\beta=a\), \(p=q\), \(H=2^a\), and support radius \(R=1\).
For \(0\le j\le q-2\), the parameters in
\eqref{h:local:parameters} give
\begin{equation}\label{h:proj:margin}
 a-j-m\ge n/2-2q+2\ge n/4+2,
 \qquad \frac{m}{a-j-m}\le\frac12.
\end{equation}
Thus \(J_j\le\exp(Cd)\), and the stated norm bound follows.
\end{proof}

\subsection{The common projected support}

A density bound alone cannot control a bounded routing function over
all label space. We also need a common support of controlled average
volume. The exact image ellipsoid supplies it.

\begin{lemma}[Exact projected volume, including short blocks]
\label{h:ev:ellipsoid}
Let \(1\le m\le d\), \(z\in\mathbb R^d\), and \(R>0\).  Use the Borel set \(\mathcal E_{z,R}\) from
\eqref{h:dv:support-sets}.  There is an absolute
\(C_{\mathrm E}\ge1\) such that, for a standard Gaussian
\(m\)-by-\(d\) matrix,
\begin{align}
 \lambda_m(\mathcal E_{z,R})
 &=v_mR^m\mathbb E_A\sqrt{\det(AA^{\mathsf T})}
 \notag\\
 &\le v_mR^m d^{m/2}
 \le R^m(2\pi e d/m)^{m/2}
 \le C_{\mathrm E}^dR^m .
 \label{h:ev:ellipsoid-volume}
\end{align}
Every density of the exact image of a measure supported on \(B(z,R)\)
vanishes almost everywhere outside \(\mathcal E_{z,R}\).
\end{lemma}

\begin{proof}
The image measure is carried by \(\mathcal E_{z,R}\), which proves
the support assertion.

Almost surely \(A\) has full row rank.  In singular-value coordinates,
\(A B(0,R)\) is an ellipsoid whose semiaxes are \(R\) times the
\(m\) singular values of \(A\).  Its volume is therefore exactly
\(v_mR^m\sqrt{\det(AA^{\mathsf T})}\).  Hadamard's inequality bounds
the square root by the product of the row norms: Gram--Schmidt replaces
each row norm by the length of its component perpendicular to the
preceding rows, which can only decrease it.  Those rows are
independent and each has expected norm at most
\(\sqrt{\mathbb E\|a_i\|^2}=\sqrt d\), giving the first inequality.
Also
\[
 (2\pi)^{m/2}
 =\int_{\mathbb R^m}e^{-\|y\|^2/2}\,dy
 \ge e^{-m/2}v_m m^{m/2},
\]
which proves \(v_m\le(2\pi e/m)^{m/2}\).

To see uniformity for short blocks, put \(t=m/d\in(0,1]\).  Since
\(t\log(1/t)\le1/e\),
\[
 \frac1d\log\!\left((2\pi e d/m)^{m/2}\right)
 =\frac t2\log(2\pi e)+\frac t2\log(1/t)
 \le\frac12\log(2\pi e)+\frac1{2e}.
\]
This supplies an absolute \(C_{\mathrm E}\) in
\eqref{h:ev:ellipsoid-volume}.  For any full-block choice
\(m=p\) with \(d/16\le p\le d\), one can instead use \(d/p\le16\)
directly to obtain
\begin{equation}\label{h:ev:full-ellipsoid}
 \lambda_p(\mathcal E_{z,R})
 \le \bigl(\sqrt{32\pi e}\,R\bigr)^p.
\end{equation}
\end{proof}
